\BeleskeID{02-s011}
% element: 02-s011:e02
% element: 02-s011:e03
% element: 02-s011:d01
Prvi korak je dodela logičkih promenljivih ulaznim i izlaznom digitalnom signalu. Ulaze označavamo sa \(C\), \(B\) i \(A\), a izlaz sa \(F\). U pozitivnoj logici, kada je fizički ulazni signal na nivou logičke jedinice, odgovarajuća promenljiva ima vrednost 1; nivou logičke nule odgovara vrednost 0. Isto povezivanje primenjujemo i na izlaz.

Na osnovu zahteva postavljenih pred sistem formiramo funkcionalnu tabelu. Tri binarna ulaza daju \(2^3=8\) kombinacija. Za svaki red posmatramo broj ulaza na nivou logičke jedinice i prema zahtevu dodeljujemo vrednost \(F\).

Uslov „dva od tri ulaza“ znači \emph{tačno dva}, a ne „najmanje dva“. Jedinice dobijamo za 011, 101 i 110. Kada su sva tri ulaza 1, uslov nije ispunjen i izlaz ostaje 0. To razlikuje zadatak od većinske funkcije, koja bi i za 111 dala 1. Vrednosti izlaza zato ne određujemo prema numeričkoj veličini binarnog broja, već prema postavljenom uslovu.
